\documentclass[serif]{beamer}
\input{sns_beamer_preamble.tex}

\begin{document}

\begin{frame}{Real Induction}
    우리가 고등학교에서 처음 접하는 수학적 귀납법은 자연수에 대한 보편 양화 명제의 증명에 \\
    이용하는 명제입니다.
    \begin{hint}[수학적 귀납법]
        자연수 전체의 집합 $\mathbb N$의 부분집합 $S$에 대해 다음이 성립한다.
        \[
            [1\in S~\wedge~\forall n\, (n\in S~\Longrightarrow~n+1\in S)]~\Longrightarrow~S=\mathbb N
        \]
    \end{hint}
    실수 집합 $\mathbb R$에서도 비슷한 사고방식을 적용할 수 있습니다. \textcolor{blue}{Pete L. Clark}의
    \begin{center}
        \textcolor{blue}{The Instructor's Guide to Real Induction}
    \end{center}
    을 통해 $\mathbb R$ 위의 귀납법, Real Induction을 살펴봅시다. [arXiv:1208.0973]
    \begin{dfn*}{}
        두 실수 $a<b$에 대하여, 닫힌구간 $[a,b]$의 부분집합 $S$가 다음을 만족하면 \\
        \textbf{귀납적(inductive)}이라 한다.\\[3pt]
        (RI1) $a\in S$. \\
        (RI2) $a\leq x<b$이고 $x\in S$이면, 어떤 $y>x$에 대하여 $[x,y]\subseteq S$이다. \\
        (RI3) $a<x\leq b$이고 $[a,x)\subseteq S$이면, $x\in S$이다.
    \end{dfn*}
    (RI1)은 수학적 귀납법의 base case, (RI2)와 (RI3)는 inductive step에 대응합니다. \\
    그러므로 $[a,b]$의 귀납적인 부분집합 $S$는 $[a,b]$ 전체와 같아야 함은 쉽게 예상할 수 있습니다.
\end{frame}

\begin{frame}{}
    \begin{thm*}{Real Induction}
        닫힌구간 $[a,b]$의 부분집합 $S$에 대하여 다음은 서로 동치이다.
        \begin{enumerate}[(i)]
            \item $S$는 귀납적이다.
            \item $S=[a,b]$.
        \end{enumerate}
    \end{thm*}
    \begin{proof}
        (ii) $\Rightarrow$ (i)은 자명하다. $S$가 귀납적이고 $S^\prime:=[a,b]\setminus S$가 공집합이 아님을 가정하자. \\
        그러면 $\nu:=\inf S^\prime<\infty$가 존재한다.
        \begin{enumerate}[1.]
            \item $\nu=a$라 하자. (RI1)에서 $a\in S$이므로 (RI2)에 의해 어떤 $y>a$에 대하여 $[a,y]\subseteq S$\\이다.
                이는 $y$가 $a$보다 큰 $S^\prime$의 하계임을 의미하므로 모순이다. 따라서 $\nu>a$이다.
            \item 이제 $\nu\in S$라 하자. 만약 $\nu=b$이면 $S=[a,b]$이다. $\nu<b$이면 (RI2)에 의해 어떤 $y^\prime>\nu$에 대하여 $[\nu,y^\prime]\subseteq S$이고,
                이 역시 $\nu$의 정의에 모순이다.
            \item 마지막으로 $\nu\in S^\prime$이라 하자. $[a,\nu)\subseteq S$이므로 (RI3)에 의해 $\nu\in S$이므로 모순이다.
        \end{enumerate}
        그러므로 $S^\prime$은 공집합이고, $S=[a,b]$이다.
    \end{proof}
    귀납적 집합의 정의에서 $a$에서 $b$까지 구간을 `연속적으로' 채워나가는 느낌, 그리고 위의 정리의 증명에서 하한의 존재성을 이용하는 점을 통해
    실수의 완비성과 Real Induction의 관계를 유추할 수 있습니다.
\end{frame}

\begin{frame}{}
    일반적인 전순서 집합(totally ordered set)에서 완비성과 Real Indcution의 관계를 살펴봅시다.
    전순서 집합 $X$가 다음 두 동치 조건 중 하나를 만족하면 \textbf{데데킨트 완비성(Dedekind completeness)}을 갖는다고 합니다.
    \begin{itemize}
        \item[-] 공집합이 아니고 위로 유계인 임의의 부분집합이 상한을 갖는다.
        \item[-] 공집합이 아니고 아래로 유계인 임의의 부분집합이 하한을 갖는다.
    \end{itemize}
    편의를 위해, 전순서 집합 $X$의 최대 원소를 $\top$, 최소 원소를 $\bot$라 합시다.
    논의에 앞서, 기존 정의를 약간 수정하여 전순서 집합 위의 귀납적 집합을 정의하겠습니다.
    \begin{dfn*}{}
        전순서 집합 $(X,\leq)$의 부분집합 $S$가 다음을 만족하면 \textbf{귀납적(inductive)}이라 한다. \\[3pt]
        (IS1) 어떤 $a\in X$에 대하여 $(-\infty,a]\subseteq S$이다. \\
        (IS2) 모든 $x\in S$에 대하여, $x=\top$이거나 어떤 $y>x$에 대하여 $[x,y]\subseteq S$이다. \\
        (IS3) 모든 $x\in X$에 대하여, $(-\infty, x)\subseteq S$이면 $x\in S$이다.
    \end{dfn*}
\end{frame}

\begin{frame}{}
    \begin{thm*}{}
        전순서 집합 $X$에 대하여 다음은 서로 동치이다.
        \begin{enumerate}[(i)]
            \item $X$는 데데킨트 완비성을 갖는다.
            \item $X$의 귀납적인 부분집합은 $X$뿐이다.
        \end{enumerate}
    \end{thm*}
    \begin{proof}
        \begin{itemize}
            \item[($\Rightarrow$)] 먼저 $X$가 데데킨트 완비성을 갖는다고 하고, $S\subseteq X$가 귀납적이라 하자.
                이전과 같은 방식으로, $S^\prime:=X\setminus S$가 공집합이 아니라 가정하자.
                (IS1)을 만족시키는 $a$를 하나 잡으면 $a$는 $S^\prime$의 하계가 되므로 $\nu:=\inf S^\prime$가 존재한다.
                $\nu$보다 작은 모든 원소는 $S$에 속하므로 (IS3)에 의해 $\nu\in S$이다.
                만약 $\nu=\top$이면 $S^\prime$은 $\{\top\}$이거나 $\varnothing$이고, 이는 모두 모순이다.
                즉 $\nu<\top$이고, (IS2)에 의해 어떤 $y>\nu$에 대하여 $[\nu,y]\subseteq S$이다.
                따라서 $(-\infty,y]\subseteq S$이고 $y$는 $\nu$보다 큰 $S^\prime$의 하계이므로 이는 모순이다.
            \item[($\Leftarrow$)] 이제 $X$가 데데킨트 완비성을 갖지 않는다고 하자.
                공집합이 아닌 $X$의 부분집합 $T$가 아래로 유계이지만 하한이 존재하지 않는다고 하자.
                $S$를 $T$의 모든 하계의 집합이라 하자. \\
                (IS1) $T$의 한 하계 $a$에 대하여 $(-\infty,a]\subseteq S$이다.
                (IS2) $x\in S$이고 $x\ne\top$이면, $x$는 $T$의 하한이 아니므로 $x$보다 큰 $S$의 원소 $y$가 존재한다.
                따라서 $S$의 정의에 의해 $[x,y]\subseteq S$이다.
                (IS3) 만약 $(-\infty,x)\subseteq S$이지만 $x\notin S$인 $x\in X$가 존재한다면, 어떤 $t\in T$에 대하여 $t<x$이다.
                $t\in S$이므로 $t$는 $T$의 최소 원소이자 하한이 되고, 이는 모순이다.
                따라서 $S$는 귀납적인 $X$의 진부분집합이다.\qedhere
        \end{itemize}
    \end{proof}
    위 정리에서, $X$를 정렬 집합(well-ordered set)으로 두면 초한귀납법(transfinite induction), 특히 자연수에 대한 수학적 귀납법을 얻습니다.
    $X$를 $\mathbb R$의 닫힌구간으로 \\ 두면 앞서 살펴본 Real Induction을 얻습니다.
\end{frame}

\begin{frame}{}
    이제 Real Induction으로 해석학에서 등장하는 정리를 증명하는데 이용해 봅시다.
    \begin{thm*}{사잇값 정리(Intermediate Value Theorem; IVT)}
        연속함수 $f:[a,b]\to\mathbb R$에 대하여 $k$를 $f(a)$와 $f(b)$ 사이의 수라 하면, \\
        어떤 $c\in[a,b]$에 대하여 $f(c)=k$이다.
    \end{thm*}
    \begin{proof}
        다음을 증명하자: 함수 $f:[a,b]\to\mathbb R\setminus\{0\}$가 연속이고 $f(a)>0$이면 $f(b)>0$이다.
        \[
            S=\{x\in[a,b]\mid f(x)>0\}
        \]
        이라 하고, $S=[a,b]$임을 보이자. (RI1) 가정에 의해 $f(a)>0$이므로 $a\in S$이다. \\ (RI2) $x\in S$, $x<b$라 하자.
        함수 $f$가 $x$에서 연속이므로 어떤 $\delta>0$에 대해 함수 $f$는 구간 $[x,x+\delta]$에서 양의 값을 갖고, $[x,x+\delta]\subseteq S$이다.
        (RI3) $x\in(a,b]$이고 $[a,x)\subseteq S$라 하자. $f(x)>0$가 아니라면 $f(x)<0$이어야 한다. 마찬가지로 연속성에 의해
        어떤 $\delta>0$에 대해 함수 $f$는 구간 $[x-\delta,x]$에서 음의 값을 가져야 하고, 이는 모순이다.
    \end{proof}
\end{frame}

\begin{frame}{}
    \begin{thm*}{Bolzano--Weierstrass}
        닫힌구간 $[a,b]$의 임의의 무한 부분집합은 극한점을 갖는다.
    \end{thm*}
    \begin{proof}
        $\mathcal{A}\subseteq[a,b]$라 하고, 다음을 만족하는 모든 $x\in[a,b]$의 집합을 $S$라 하자.
        \begin{center}
            $\mathcal{A}\cap[a,b]$가 무한집합이면, 극한점을 갖는다.
        \end{center}
        $S=[a,b]$임을 보이자. (RI1)은 자명하다. (RI2) $x\in S\cap[a,b)$라 하자.
        \begin{enumerate}[1.]
            \item 만약 $\mathcal{A}\cap[a,x]$가 무한집합이면, 이는 극한점을 가지므로 $\mathcal{A}\cap[a,b]$ 역시 그렇다. \\
                이 경우 $S=[a,b]$이다.
            \item 어떤 $\delta>0$에 대해 $\mathcal{A}\cap[a,x+\delta]$가 유한집합이면, $[x,x+\delta]\subseteq S$이다.
            \item 아니면 모든 $\delta>0$에 대해 $\mathcal{A}\cap[a,x]$는 유한집합이지만 $\mathcal{A}\cap[a,x+\delta]$는 무한집합이다. 
                그러면 $x$가 $\mathcal A$의 극한점이 되므로 $S=[a,b]$이다.
        \end{enumerate}
        (RI3) $[a,x)\subseteq S$이면 다음 경우가 가능하다.
        \begin{enumerate}[1.]
            \item $\mathcal{A}\cap[a,y]$가 무한집합이 되도록 하는 $y<x$가 존재한다.
            \item $\mathcal{A}\cap[a,x]$가 유한집합이다.
            \item 모든 $y<x$에 대해 $\mathcal{A}\cap[a,y]$는 유한집합이고, $\mathcal{A}\cap[a,x]$는 무한집합이다.
        \end{enumerate}
        세 경우 모두 $x\in S$를 함의한다.
    \end{proof}
\end{frame}

\begin{frame}{}
    이제 $\mathbb R$의 닫힌구간의 연결성과 콤팩트성을 보겠습니다. 해석학 교재의 증명과 비교해 보세요!
    \begin{thm*}{}
        $\mathbb R$의 닫힌구간 $[a,b]$는 연결집합이다.
    \end{thm*}
    \begin{proof}
        구간 $[a,b]$의 부분집합 $A$가 열린집합이면서 닫힌집합이라 하자. $a\in A$라 가정하고 $A=[a,b]$임을 보이자. (RI1)은 가정에 의해 성립한다.
        $A$는 열린집합이므로 (RI2)가 \\ 성립하고, $A$는 닫힌집합이기도 하므로 (RI3) 역시 성립한다.
    \end{proof}
    \begin{thm*}{Heine--Borel}
        $\mathbb R$의 닫힌구간 $[a,b]$는 콤팩트집합이다.
    \end{thm*}
    \begin{proof}
        $\{U_i\}_{i\in I}$를 구간 $[a,b]$의 열린 덮개라 하고,
        \[
            S=\{x\in[a,b]\mid \{U_i\cap[a,x]\}_{i\in I}\text{가 유한 부분덮개를 갖는다.}\}
        \]
        로 정의하자. $S=[a,b]$임을 보이자. (RI1)은 자명하다.
        (RI2) $U_1,\cdots,U_n$이 $[a,x]$의 유한 부분덮개이면 어떤 $U_k$는 $\delta>0$에 대해 구간 $[x,x+\delta]$를 포함한다.
        (RI3) $[a,x)\subseteq S$, $x\in U_{i_x}$라 하자. $[x-\delta,x]\subseteq U_{i_x}$인  $\delta>0$를 잡자.
        $x-\delta\in S$이므로 $[a,x-\delta]$의 유한 부분덮개 $U_1,\cdots U_n$이 존재한다. 따라서 $\{U_k\}_{k=1}^n\cup\{U_{i_x}\}$는 $[a,x]$의 유한 부분덮개이다.
    \end{proof}
\end{frame}

\begin{frame}{}
    Real Induction과 마찬가지로, 전순서 집합의 연결성과 콤팩트성을 논할 수 있습니다. 먼저 연결성은 조밀성이 필요합니다.
    임의의 두 원소 $a<b$에 대해 $a<c<b$인 원소 $c$가 \\ 존재하는 전순서 집합을 \textbf{조밀하다(dense)}고 합니다.
    \begin{lemma*}{}
        데데킨트 완비성을 갖는 전순서 집합 $(X,\leq)$의 부분집합 $S$가 $X$의 순서 위상에서 \\
        닫힌집합이면, $(S,\leq)$도 데데킨트 완비성을 갖는다.
    \end{lemma*}
    \begin{proof}
        정의를 이용하여 쉽게 보일 수 있으므로 생략한다.
    \end{proof}
    \begin{thm*}{}
        전순서 집합 $X$에 대하여 다음은 서로 동치이다.
        \begin{enumerate}[(i)]
            \item $X$는 데데킨트 완비성과 조밀성을 갖는다.
            \item $X$는 순서 위상에서 연결성을 갖는다.
        \end{enumerate}
    \end{thm*}
    \begin{proof}
        \begin{itemize}
            \item[($\Rightarrow$)] 먼저 $X$에 최소 원소 $\bot$가 존재한다고 하자. $S\subseteq X$가 $\bot$을 갖는 열린집합이자 닫힌집합이면 $S$는 귀납적이다.
                그러므로 $X$의 데데킨트 완비성에 의해 $S=X$이고, $X$는 연결성을 갖는다. 이제 $X\ne\varnothing$인 일반적인 경우를 보자. $a\in X$라 하면 보조정리로부터
                $[a,\infty)$와 $(-\infty,a]$ 모두 연결집합임을 알 수 있다. $X=(-\infty,a]\cup[a,-\infty)$이므로 $X$는 연결성을 갖는다.
            \item[($\Leftarrow$)] 정의를 이용하여 쉽게 보일 수 있으므로 생략한다.
        \end{itemize}
    \end{proof}
\end{frame}

\begin{frame}{}
    전순서 집합 $X$의 모든 부분집합이 상한 (또는 하한)을 가지면 $X$가 \textbf{완비성(completeness)}을 갖는다고 합니다. (데데킨트 완비성과 다름에 유의하세요!)
    \begin{thm*}{}
        전순서 집합 $X$에 대하여 다음은 서로 동치이다.
        \begin{enumerate}[(i)]
            \item $X$는 완비성을 갖는다.
            \item $X$는 순서 위상에서 콤팩트성을 갖는다.
        \end{enumerate}
    \end{thm*}
    \begin{proof}
        \begin{itemize}
            \item[($\Rightarrow$)] $\{U_i\}_{i\in I}$를 $X$의 열린 덮개라 하자. 다음을 만족하는 원소 $x\in X$의 집합을 $S$라 하자.
                \begin{center}
                    $[\bot,x]$의 열린 덮개 $\{U_i\cap[0,x]\}_{i\in I}$는 유한 부분덮개를 갖는다.
                \end{center}
                (IS1) $\bot\in S$임은 자명하다. (IS2) $U_1\cap[0,x], U_2\cap[0,x], \cdots, U_n\cap[0,x]$이 $[0,x]$의 유한 부분덮개라 하자. 만약 $[x,y]=\{x,y\}$가 되도록 하는 원소 $y\in X$가 존재하면 $y$를 포함하는 $U_y$를 포함시키면 된다. 그렇지 않다면, $x$를 포함하는 어떤 $U_x$는 $y>x$에 대하여 닫힌구간 $[x,y]$를 포함한다. (IS3) 원소 $x\in X$에 대하여 $[\bot,x)\in S$라 하고, $x$를 포함하는 $U_x$를 택하자. $U_x$는 어떤 $y<x$에 대하여 닫힌구간 $[y,x]$를 포함하므로 정의에 의해 $x\in S$이다. $X$는 데데킨트 완비성을 갖고, $S$는 귀납적이므로 $S=X$이다. 즉 $\top\in S$이므로 $\{U_i\}_{i\in I}$는 유한 부분덮개를 갖는다.
            \item[($\Leftarrow$)] 생략. (직접 해보세요!) \qedhere
        \end{itemize}
    \end{proof}
\end{frame}

\begin{frame}{}
    데데킨트 완비성에 주목하면 다음 따름정리를 얻습니다.
    \begin{cor*}{일반화된 Heine--Borel}
        전순서 집합 $X$에 대하여 다음은 서로 동치이다.
        \begin{enumerate}[(i)]
            \item $X$는 데데킨트 완비성을 갖는다.
            \item 순서 위상에서 $X$의 부분공간 $S$가 콤팩트일 필요충분조건은 \\ $S$가 유계 닫힌집합인 것이다.
        \end{enumerate}
    \end{cor*}
    순서체에서는 지금까지의 논의가 다음 따름정리로 요약됩니다.
    \begin{cor*}{}
        순서체 $F$에 대하여 다음은 모두 동치이다.
        \begin{enumerate}[(i)]
            \item $F\simeq\mathbb R$.
            \item $F$의 임의의 두 원소 $a<b$에 대하여, $S\subseteq[a,b]$가 (IS1), (IS2), (IS3)을 \\ 모두 만족하면 $S=[a,b]$이다.
            \item $F$의 모든 닫힌구간 $[a,b]$는 순서 위상에서 연결집합이다.
            \item $F$의 모든 닫힌구간 $[a,b]$는 순서 위상에서 콤팩트집합이다.
        \end{enumerate}
    \end{cor*}
\end{frame}
\end{document}